Throughout, $B$ is the polar form of $Q$, and blades $e_I f_K$ are the increasing products in the basis $e_1, \dots, e_m, f_1, \dots, f_r$ of $V$; by \claimref{std-algebra.blade-basis} they form a basis of $\Cl(V, Q)$, degenerate forms included.

(i) For $f \in \operatorname{rad} V$ and $v \in V$, \claimref{std-algebra.vector-squares} gives $fv + vf = 2B(f, v) = 0$, so $f v = -v f$. For a product $a = v_1 \cdots v_k$ of vectors this gives $f a = (-1)^{k} a f = \hat{a} f$, and by linearity $f a = \hat{a} f$ for every $a$, since products of vectors span the algebra. Hence $a f c = a \hat{c} f \in \Cl\, \operatorname{rad} V$, so the two-sided ideal $J$ equals the left ideal $\Cl \operatorname{rad} V$, and by the same identity read the other way it equals the right ideal $\operatorname{rad} V\, \Cl$. Also $f^{2} = Q(f) = B(f, f) = 0$ and $f_i f_j = -f_j f_i$. A spanning set of $\Cl \operatorname{rad} V$ is $\{e_I f_K f_j\}$, and $e_I f_K f_j$ is $0$ if $j \in K$ and $\pm e_I f_{K \cup \{j\}}$ otherwise. So $J \subseteq \operatorname{span}\{e_I f_K : K \ne \emptyset\}$; conversely, if $K \ne \emptyset$ with largest index $k$, then $e_I f_K = (e_I f_{K \setminus \{k\}}) f_k \in J$. The blades with $K \ne \emptyset$ are independent, and there are $2^{n} - 2^{m}$ of them.

(ii) By (i), $J^{k}$ is spanned by products $a_1 f_{j_1} a_2 f_{j_2} \cdots a_k f_{j_k}$, and moving every radical factor to the right with $f a = \hat a f$ turns such a product into $\pm a_1 \hat a_2 \cdots f_{j_1} \cdots f_{j_k}$. The radical factors anticommute and square to zero, so the product of $k$ of them is $\pm f_K$ with $\lvert K \rvert = k$, or $0$. Hence $J^{k} \subseteq \Cl \cdot \operatorname{span}\{f_K : \lvert K \rvert = k\}$, which is the span of the blades $e_I f_L$ with $\lvert L \rvert \ge k$. Conversely such a blade is $(e_I f_{L'} f_{l_1}) f_{l_2} \cdots f_{l_k}$ with $l_1 < \dots < l_k$ the last $k$ indices of $L$, a product of $k$ elements of $J$. For $k = r + 1$ there is no such blade, so $J^{r+1} = 0$; for $k = r$ the blades are $e_I f_1 \cdots f_r$, nonzero by the blade basis, so the index is exactly $r + 1$ when $r \ge 1$.

(iii) $Q(v + f) = Q(v) + 2B(v, f) + Q(f) = Q(v)$ for $f \in \operatorname{rad} V$, so $\bar Q$ is well defined; its polar form is $\bar B(\bar u, \bar v) = B(u, v)$, and if $\bar u$ is orthogonal to all of $\bar V$ then $u \in \operatorname{rad} V$, so $\bar u = 0$: $\bar Q$ is nondegenerate. The map $v \mapsto \bar v \in \Cl(\bar V, \bar Q)$ satisfies $\bar v^{2} = \bar Q(\bar v) = Q(v)$, so the universal property (\claimref{std-algebra.existence-uniqueness}) gives a unique algebra homomorphism $\pi$ with $\pi(v) = \bar v$. It is onto since $\bar V$ generates. It kills $\operatorname{rad} V$, hence $J \subseteq \ker \pi$. The restriction $U \to \bar V$ is a linear bijection with $\bar Q(\bar u) = Q(u)$, an isometry of $(U, Q|_U)$ onto $(\bar V, \bar Q)$; so $(U, Q|_U)$ is nondegenerate, $\Cl(U, Q|_U)$ is the subalgebra of $\Cl(V, Q)$ spanned by the $e_I$ (\claimref{std-algebra.functoriality}), and $\pi$ maps it isomorphically onto $\Cl(\bar V, \bar Q)$ (\claimref{std-algebra.functoriality} again, both sides of dimension $2^{m}$). Counting, $\dim \Cl(V, Q)/J = 2^{n} - (2^{n} - 2^{m}) = 2^{m} = \dim \Cl(\bar V, \bar Q)$, so $\ker \pi = J$, and $\Cl(V, Q) = \Cl(U, Q|_U) \oplus J$ since the $e_I$ and the $e_I f_K$ ($K \ne \emptyset$) together form the blade basis.

(iv) Choose the basis of $U$ orthogonal with $Q(e_i) = \pm 1$, possible since $(U, Q|_U)$ is nondegenerate (\claimref{std-algebra.sylvester-inertia}). Left multiplication $L_x$ by a blade $x = e_I f_K \ne 1$ sends each blade to a multiple of a blade with a different index set (the symmetric difference with $I \cup K$), so its diagonal in the blade basis is zero; hence $\operatorname{tr} L_x = 2^{n} \langle x \rangle_0$ for every $x$. If $x$ is nilpotent so is $L_x$, and $\operatorname{tr} L_x = 0$, so $\langle x \rangle_0 = 0$. Let $N$ be a left ideal with $N^{k} = 0$. For $x \in N$ and any $y$, $yx \in N$ is nilpotent, so $\langle y x \rangle_0 = 0$ for all $y$. Take $y = e_I^{-1} = \pm e_I$ for $K = \emptyset$: since $e_I e_{I'} f_{K'}$ has scalar part $0$ unless $I' = I$ and $K' = \emptyset$, this reads off the coefficient of $e_I$ in $x$, which is therefore $0$; so $x \in J$ by (i). For right ideals use $\langle x y \rangle_0$ and $y = e_I^{-1}$ on the other side. Thus every nilpotent one-sided ideal lies in $J$, and $J$ is nilpotent by (ii). The same computation shows the radical of $(x, y) \mapsto \langle xy \rangle_0$ is exactly the span of the $e_I f_K$ with $K \ne \emptyset$, that is $J$: the blades pair diagonally, $\langle e_I f_K e_I f_K \rangle_0 = \pm\prod_{i \in I} Q(e_i) \prod_{k \in K} Q(f_k)$, which is $\pm 1$ for $K = \emptyset$ and $0$ otherwise. This description involves no choices. Applied with $r = 0$ it shows that $\Cl(\bar V, \bar Q)$ has no nonzero nilpotent one-sided ideal.

For $r = 1$, (ii) gives $J^{2} = 0$ and (i) gives $\dim J = 2^{n-1}$.

Example. Let $k \in \R^{1,3}$ be null and $U = k^{\perp}$, a three-dimensional subspace containing $k$. Every $a \in U$ has $B(a, k) = 0$, so $k \in \operatorname{rad} U$; choosing $e_1, e_2 \in U$ with $U = \R k \oplus \operatorname{span}(e_1, e_2)$, the form on $\operatorname{span}(e_1, e_2)$ is negative definite (a spacelike plane orthogonal to $k$), so $\operatorname{rad} U = \R k$ and $r = 1$. The inclusion $U \to \R^{1,3}$ extends to an injective homomorphism $\Cl(U) \to \Cl(1,3)$ (\claimref{std-algebra.functoriality}), so $\Cl(U)$ is the subalgebra generated by $U$. By (i) and (ii), $J = k\,\Cl(U) = \operatorname{span}\{k, k e_1, k e_2, k e_1 e_2\}$ with $J^{2} = 0$; its bivectors $N = k a$, $a \in \operatorname{span}(e_1, e_2)$, are the generators of the null rotations $1 + N$ about $k$: any two multiply to zero, so they commute, $N^{2} = 0$, and $e^{N} = 1 + N$ (\claimref{std-algebra.bivector-exponential}, null case). By (iii), $\Cl(U)/J \cong \Cl(\operatorname{span}(e_1, e_2)) \cong \Cl(0,2)$.
